\section{Introduction}
\label{sec:intro}

\definecolor{mygreen}{HTML}{008000}
\begin{figure}[!t]
  \centering
  % \includegraphics[width=1\linewidth]{figure/fig1_ver3.pdf}
  % \includegraphics[width=1\columnwidth, trim=0cm 0cm 0cm 0cm, clip]{figure/fig1_ver4.pdf}
  \includegraphics[width=1\columnwidth, trim=0cm 0cm 0cm 0cm, clip]{figure/fig1_ver5.pdf}
  % \fbox{\rule{0pt}{1.0\columnwidth} \rule{0.9\columnwidth}{0pt}}
  \caption{Object-centric \ac{NBV} for manipulation in cluttered scenes. Instead of reconstructing the entire scene, we prioritize task-relevant regions surrounding the target object (tennis ball (\textcolor{red}{red}) or cup (\textcolor{blue}{blue})). Using these selected views with our object-aware \ac{3DGS} yields improved novel view synthesis and depth reconstruction, as shown below.}
  % cluttered scene에서 manipulation을 할 때 전체 scene을 reconsturction 하는 것보다 target object 위주 (tennis ball for green, cup for blue) 를 우선적으로 reconstruct하는 것이 효율적이다. 우리는 이를 위해 물체를 구분할 수 있는 semantic 3DGS와 거기서 얻을 수 있는 information을 기반으로 Next best view를 결정하는 policy를 사용한다.
  \label{fig:main}
  \vspace{-6mm}
\end{figure}

Accurate perception of the operating environment including the geometry of objects is essential for successful robotic manipulation. While many prior works \cite{fang2020graspnet, wu2024economic} estimate grasping points from a single view depth image, real-world cluttered scenes often contain stacked or irregularly distributed objects, where severe occlusions occur. In such cases, predefined viewpoints may either lack necessary information or include redundant observations, making them insufficient and ineffective for reliable manipulation \cite{murali20206, newbury2023deep}.  This motivates the need for a robot to \textit{actively} move and select informative viewpoints based on the scene context or the target object and incrementally refine its understanding of the scene from the newly captured perception.

% 로봇이 manipulation task를 제대로 수행하기 위해선 잡고자 하는 물체의 geometry를 포함한 surrounding environment를 정확하게 인지해야 한다. 기존에 보통 Bird's eye view에서의 single view depth image를 통해 물체의 grasping point를 얻는 연구들이 많이 진행되어 왔지만, 실제 practical한 usage에서 로봇은 여러 물체가 쌓여있거나 불규칙적으로 분포되어 있는 cluttered scene을 마주하게 된다. 이런 상황에서는 물체들간의 occlusion이 크게 작용하기 때문에 소수의 pre-define된 위치의 view만으로는 manipulation을 제대로 수행할 수 없다. 따라서 로봇에게는 scene의 상황이나 목표로 하는 물체에 따라 더 많은 정보를 줄 수 있는 view를 선택하고 그를 이용하여 점진적으로 scene을 이해해나갈 수 있는 능력이 필요하다.

Identifying a minimal set of training views that preserves reconstruction quality frames the problem of view selection and \acf{NBV} planning. Voxel-based methods \cite{zhang2023affordance, jin2025activegs} estimate volumetric completeness via per-voxel occupancy, but require high memory at finer resolution and do not account for photometric quality. Bayes’ Rays \cite{goli2024bayes} addresses this limitation by parameterizing perturbations to \ac{NeRF} \cite{mildenhall2021nerf} parameters to quantify uncertainty, yet its slow rendering speed makes it unsuitable for robotic tasks. 

% Identifying a minimal set of training views that preserves reconstruction quality frames the problem of view selection and \ac{NBV} planning. Using such an algorithm, the information gain of multiple candidate views can be quantified, and the view with the highest gain can be added to the training set to progressively complete the scene. Methods based on voxel representations \cite{zhang2023affordance, jin2025activegs, nagami2025vista} typically estimate volumetric completeness by measuring per-voxel occupancy, but this approach suffers from high memory usage depending on voxel resolution and does not account for photometric quality. Bayes’ Rays \cite{goli2024bayes} addresses this limitation by parameterizing perturbations to \ac{NeRF} \cite{mildenhall2021nerf} parameters to quantify uncertainty, yet its slow rendering speed makes it unsuitable for robotic tasks. 

% view selection과 Next best view planning은 reconstruction quality를 해치지 않는 최소한의 training view set을 찾는 문제로, 이 알고리즘을 이용해 여러 candidate view 들의 information gain을 quantify하고 그것이 가장 큰것을 새로운 training view로 추가하여 scene을 완성해나갈 수 있다. Voxel-based representation을 사용한 연구들은 voxel별로 occupancy 등을 이용하여 volumetric completeness를 비교하여 view를 선택하지만, voxel의 해상도에 따라 메모리 사용량이 크고 photometric quality를 고려하지 않는다. NeRF parameter의 perturbation을 파라미터화하여 uncertainty를 quantify한 Bayes' Rays는 이를 해소하지만 NeRF의 느린 rendering 속도로 인해 로봇 task에 적합하지 않다.

Compared to these approaches, methods \cite{jiang2024fisherrf, wilson2025pop, strong2024next} based on \acf{3DGS} \cite{kerbl20233d} offer the advantages of using an explicit representation, which enables per-Gaussian uncertainty estimation and faster rendering. By leveraging the Jacobian of Gaussian parameters with respect to the rendering output, these methods avoid heuristic modelings and select views in a direct and efficient manner. However, since well-observed Gaussians with low uncertainty still dominate the rendering contribution, these \ac{NBV} strategies tend to focus on exploitation rather than exploration. Moreover, these methods are designed for reconstructing the entire scene, making them less robust for manipulation tasks targeting objects that are occluded, small, or far from the center. In such cluttered environments, robots often require object-centric \ac{NBV} planning that can focus on a specific object.

% 문장 김.

% 이에 비해 3DGS-based method들은 explicit한 representation를 가져 per-Gaussian uncertainty를 구할 수 있고 rendering이 매우 빠르다는 장점을 가지고 있다. gaussian의 parameter들의 rendering에 대한 jacobian을 uncertainty qunatification에 사용하는 이 방법들은 heurisitc하지도 않고 view를 직접적이고 효율적으로 선택하지만, observation이 많아 certain한 Gaussian들이 여전히 rendering에 대한 기여도를 dominant하기 때문에 exploration 보다 exploitation에 집중하는 경향이 있다. 또한, 이 방법들은 모두 전체 scene에 대한 reconstruction을 목표로 할 뿐, 특정 물체에 집중된 NBV가 불가능하여 cluttered scene에서 가려지거나 작은 물체들에 대해서 manipulation을 수행할 때 robust 하지 않다.


To tackle these limitations, we propose an informative, object-centric \ac{NBV} policy integrated with object-aware \ac{3DGS}.
Our key insight is to inject instance-aware confidence score into uncertainty quantification so that information gain and view selection prioritize uncertain, task-relevant regions. This enables \ac{NBV} for both whole scene coverage and target-conditioned inspection with better novel view synthesis and depth reconstruction as shown in \figref{fig:main}.
Our object-aware \ac{3DGS} distills object instance masks from image foundation models \cite{liu2024grounding, kirillov2023segment, ren2024grounded} into a one-hot object vector, reconstructs scenes in fewer iterations than optimization using only RGB images and yields explicit per-Gaussian object features that guide \ac{NBV}.
We validate our approach through substantial improvements in both novel view synthesis and depth estimation, as well as successful deployment in real-world robotic arm experiments.
Our main contributions are as follows:

% We define a confidence score for each Gaussian, representing how likely it belongs to a specific object, and integrate this into existing uncertainty quantification to compute a semantic-aware information gain. This enables the NBV policy not only to cover previously underexplored regions of the entire scene, but also to focus viewpoint selection on specific target objects. By leveraging object masks obtained from image foundation models \cite{liu2024grounding, kirillov2023segment, ren2024grounded}, the robot’s scene representation quickly converges toward reality, and reconstruction can be achieved in fewer iterations compared to RGB-only optimization. We validate our approach through substantial performance improvements in both novel view synthesis and depth estimation, as well as successful deployment in real-world robotic arm experiments.
% 우리는 각 gaussian이 특정 물체에 얼마나 confident하게 속하는 지에 대한 score를 정의한 뒤 기존의 uncertainty quantification에 포함하여 semantic 정보를 포함하는 information gain을 계산한다. 이 방법을 통해 전체 scene을 볼 때 이제까지 제대로 보지 못한 부분을 보도록 NBV를 할 뿐만 아니라 특정 target object에 집중된 view selection을 가능하게 한다. image foundation model들을 통해 practically 얻을 수 있는 object mask를 통해 robot이 가지고 있는 scene representation은 빠르게 실제와 가까워지고, reconstruction을 반복함에도 rgb만 사용하여 optimize를 할 때보다 더 적은 iteration으로 씬을 복원할 수 있다.

% We leverage object masks from image foundation models \cite{liu2024grounding, kirillov2023segment, ren2024grounded} to supervise an additional feature for each Gaussian. 
% The resulting one-hot object vector accelerates Gaussian optimization, enables object-wise segmentation, and reweights information gain toward under-observed regions, thereby facilitating more exploration-oriented \ac{NBV}.
% 이게 어떻게 한 문장 ㅋㅋ
% The resulting one-hot object vector, which can be interpreted as an object probability, not only accelerates Gaussian optimization and enables object-wise segmentation, but also emphasizes the information gain of under-observed regions, thereby facilitating more exploration-oriented \ac{NBV} planning. 
% Furthermore, simple masking allows \ac{NBV} and reconstruction focused on specific target objects, ensuring effective and robust manipulation in cluttered scenes. 
% We validate our approach through substantial performance improvements in both novel view synthesis and depth estimation, as well as successful deployment in real-world robotic arm experiments.

% 이 문제를 해결하기 위해서 우리는 image foundation model로부터 얻을 수 있는 view들의 object mask를 이용해 Gaussian들에 추가적인 feature를 학습시켜 사용하였다. object probability로 사용될 수 있는 우리의 one-hot object vector는 Gaussian들의 학습을 가속화하고 물체별로 segment해줄 수 있을 뿐만 아니라 덜 관찰된 부분에 대한 information gain을 강조하여 NBV를 할 수 있게 해준다. 또한, 간단한 masking을 통해 특정 물체에 집중된 NBV와 reconstruction을 가능케 하여 scene이나 object에 robust한 manipulation을 보장한다. 우리는 이를 novel view synthesis 뿐만 아니라 depth estimation에서의 월등한 성능 향상과 실제 robot arm experiment에의 적용을 통해 방법론이 유효함을 확인하였다.

\begin{itemize}
    \item \textbf{Instance-aware Information Gain for Next Best View: }
    By weighting each Gaussian’s information gain with a confidence score from object features, our \ac{NBV} planner prioritizes poorly observed yet task-relevant regions, encouraging exploration.
    % object one-hot vector를 confidence로 활용하여 training view에서 rendering quality에 대한 기여도가 작은 gaussian들을 더 집중할 수 있는 information gain 식을 제시한다.
    \item \textbf{Object-aware 3D Gaussian Splatting and Fast Scene Refinement: } 
    We introduce an object-aware \ac{3DGS} that reconstructs and segments cluttered scenes in just a few iterations with low training time, by fusing incoming RGB images and object instance masks through a one-hot object vector.
    % 지속적으로 view가 추가되는 tabletop cluttered scene setting에서 rgb와 object mask를 이용해 빠르게 3d reconstruction 및 segmentation을 수행하는 sparse-view 3dgs framework을 소개한다.
    \item \textbf{Object-centric Reconstruction and Manipulation: }
    Beyond whole-scene modeling, our system can focus \ac{NBV} around user-specified objects and produces optimal 6-DoF grasp poses in cluttered scenes.
    % 전체 3d scene뿐만 아니라 robotic manipulation 시 practical한 활용을 위해 목표로 하는 물체 주변에 집중하고 최적의 6dof grasp를 결정하도록 한다.
\end{itemize}

